\documentclass[11pt, a4paper]{article}

\usepackage[utf8]{inputenc}
\usepackage[T1]{fontenc}
\usepackage{amsmath, amssymb, amsfonts}
\usepackage{graphicx}
\usepackage[margin=1in]{geometry}
\usepackage{hyperref}
\usepackage[numbers, sort&compress]{natbib}
\usepackage{booktabs}
\usepackage{siunitx}
\usepackage{caption}
\usepackage{float}
\usepackage{xcolor}

\title{\textbf{Nonlinear-QND Limits and Hardware-Realistic Optimization of Dispersive Readout Pulses in cQED}}
\author{GitHub Copilot \\ cQED Autonomous Research Platform}
\date{\today}

\begin{document}

\maketitle

\begin{abstract}
We revisit dispersive readout-pulse optimization after the main simplifying weakness of the
previous study is removed: the replay stack now includes both hardware-realistic waveform
transport and an explicit family-dependent strong-drive disturbance channel. The study is
implemented entirely on top of \texttt{cqed\_sim}. Native hardware replay is handled with
\texttt{SequenceCompiler(hardware=...)} and a study-local effective mixing layer maps
occupancy- and slew-dependent readout backaction onto explicit multilevel ancilla
transition drives. We compare square, smoothed-square, ring-up/hold/ring-down,
free procedural segments, analytic nulling tails, Fourier-basis pulses, and a higher
dimensional piecewise reference.

The main answer is qualified. Procedural structure remains valuable, but the freest pulse
family is no longer the best practical default once nonlinear-QND penalties and transport
distortion are included. Across the executed frontier at \SIlist{96;240;496}{ns}, the best
rich-model family is the low-dimensional \texttt{ring\_hold} design, with detector-limited
matched-filter fidelities $0.6480$, $0.9599$, and $0.9982$, respectively. At the
representative \SI{240}{ns} point, \texttt{procedural\_segments} still attains very high
information fidelity ($F_\eta = 0.9547$), but it does so with a larger measurement-induced
transition probability ($4.66\times10^{-3}$) and lower QND preservation ($0.9886$) than
smoother structured alternatives. Hardware replay accounts for most of the nominal fidelity
loss, while the nonlinear layer accounts for most of the new QND penalty. A higher
dimensional reference improves the information-only objective but is not competitive under
balanced practical scoring.

The resulting recommendation is pragmatic: low-dimensional physics-informed pulses remain
near-optimal, but unrestricted procedural freedom is not the right default under a
QND-aware, hardware-realistic model. A ring-up/hold/ring-down pulse is the best default
family in the present study, while free procedural pulses should be reserved for cases
where the user can tolerate a tighter QND budget and re-optimize against the actual
hardware stack.
\end{abstract}

\section{Introduction}
Dispersive readout is the dominant measurement mechanism for superconducting qubits in
circuit QED because it converts the qubit state into a state-conditioned cavity response
that can be discriminated with low added noise \cite{blais2021circuit,koch2007transmon}.
That basic picture is familiar. The engineering difficulty is that useful readout must
also empty the cavity quickly, preserve repeated-use readiness, and avoid drive-induced
state disturbance \cite{jeffrey2014fast,mcclure2016rapid,sank2016measurement}.

The predecessor study in this repository showed that low-dimensional procedural readout
pulses can outperform square baselines in a conventional dispersive replay stack. The open
question was whether that conclusion survives once two missing ingredients are restored:
hardware-realistic waveform transport and a family-dependent non-QND mechanism that turns
on under strong readout occupancy. This follow-up study addresses exactly that question.

The answer matters because it separates two different claims. One claim is that control
structure helps; the earlier study already supported that. The stronger claim is that the
freest procedural family is still the best practical choice once the readout stack is made
more realistic. The present work shows that the first claim survives while the second does
not survive cleanly.

The remainder of the paper is organized as follows. Section 2 defines the model, hardware
assumptions, and optimization protocol. Section 3 reports the frontier comparison,
representative rich-model diagnostics, transport effects, and robustness results.
Section 4 documents the validation checks. Section 5 interprets the control conclusion,
and Sections 6 and 7 give the final recommendation and handoff limitations.

\section{System and Methods}

\subsection{Hamiltonian}
The base replay model is the multilevel dispersive transmon--readout Hamiltonian supported
by \texttt{cqed\_sim},
\begin{equation}
\frac{H_0}{\hbar} =
\delta_c a^\dagger a +
\delta_q b^\dagger b +
\frac{\alpha}{2} b^{\dagger 2} b^2 +
\frac{K}{2} a^{\dagger 2} a^2 +
\chi\, a^\dagger a\, b^\dagger b,
\end{equation}
where $a$ is the readout-mode annihilation operator, $b$ is the transmon lowering
operator, $\alpha$ is the transmon anharmonicity, $K$ is the cavity self-Kerr, and
$\chi$ is the dispersive shift. The cavity is driven through a complex envelope
$\epsilon_c(t)$, and the richer study-local layer adds effective multilevel transition
drives,
\begin{equation}
\frac{H_{\mathrm{drive}}}{\hbar} =
\epsilon_c(t)a^\dagger + \epsilon_c^*(t)a +
\Omega_{ge}(t)\sigma_{ge}^\dagger + \Omega_{ge}^*(t)\sigma_{ge} +
\Omega_{ef}(t)\sigma_{ef}^\dagger + \Omega_{ef}^*(t)\sigma_{ef}.
\end{equation}
The synthetic envelopes $\Omega_{ge}(t)$ and $\Omega_{ef}(t)$ are not ad hoc simulator
replacements. They are explicit drives inside the normal \texttt{cqed\_sim} replay path,
with amplitudes determined by the distorted readout waveform, the estimated instantaneous
occupancy $n(t)/n_{\mathrm{crit}}$, and the local phase slew. In this way the richer model
can generate family-dependent induced transitions without leaving the supported pulse,
compiler, and Lindblad workflow.

Dissipation is modeled with cavity damping, thermal occupation, dephasing, and multilevel
relaxation via \texttt{NoiseSpec(kappa, tphi, transmon\_t1, nth)}. The baseline ladder
times are $(T_{1,ge}, T_{1,ef}) = (\SI{30}{\micro\second}, \SI{12}{\micro\second})$.

\subsection{Simulation Parameters}
\begin{table}[H]
  \centering
  \caption{Nominal physical, hardware, and nonlinear-model parameters.}
  \begin{tabular}{@{} l l l @{} }
    \toprule
    Parameter & Value & Unit \\
    \midrule
    $\omega_q/2\pi$ & 6.150 & GHz \\
    $\omega_r/2\pi$ & 8.597 & GHz \\
    $\alpha/2\pi$ & -255 & MHz \\
    $\chi/2\pi$ & -2.84 & MHz \\
    $\kappa/2\pi$ & 2.4 & MHz \\
    $K/2\pi$ & -28 & kHz \\
    Transmon truncation $n_{\mathrm{tr}}$ & 3 & levels \\
    Cavity truncation $n_{\mathrm{cav}}$ & 14 & Fock states \\
    $T_{1,ge}$ & 30 & $\mu$s \\
    $T_{1,ef}$ & 12 & $\mu$s \\
    $T_\phi$ & 30 & $\mu$s \\
    Thermal occupation $n_{\mathrm{th}}$ & 0.015 & --- \\
    Detector efficiency $\eta$ & 0.35 & --- \\
    Time step $\Delta t$ & 4 & ns \\
    Max cavity drive $|\epsilon_{\max}|/2\pi$ & 4.0 & MHz \\
    Estimated $n_{\mathrm{crit}}$ & 20.3 & photons \\
    Low-pass bandwidth & 75 & MHz \\
    IQ gain ratio $g_Q/g_I$ & 0.94 & --- \\
    Quadrature skew & 5 & deg \\
    Image leakage & 0.03 & --- \\
    Amplitude quantization & 9 & bits \\
    Timing quantum & 8 & ns \\
    Mixing onset ratio & 0.10 & $n/n_{\mathrm{crit}}$ \\
    $ge$ mixing scale & 0.09 & --- \\
    $ef$ mixing scale & 0.05 & --- \\
    \bottomrule
  \end{tabular}
\end{table}

\subsection{Computational Approach}
All simulations and optimizations were implemented in
\path{studies/nonlinear_qnd_hardware_realistic_readout/scripts/readout_opt/}
using supported \texttt{cqed\_sim} APIs \cite{cqedsim}. The main reusable pieces are the
family parametrizations, hardware-aware replay functions, metric construction, and study
orchestration helpers.

We compare six practical families: square, smoothed-square, ring-up/hold/ring-down,
free three-segment procedural pulses, analytic nulling tails, and a Fourier basis. A
10-segment piecewise waveform is used only as a higher-dimensional reference. The frontier
sweep was executed at \SIlist{96;240;496}{ns} and saved in
\texttt{data/frontier\_summary.json}. The deeper representative analysis at
\SI{240}{ns}, including regime breakdown, reference slices, robustness, QND stress, and
convergence probes, was saved in \texttt{data/study\_summary.json}.

Three replay hierarchies were used. The \emph{full} regime is the earlier nominal replay
with ordinary multilevel dissipation. The \emph{hardware} regime compiles the programmed
envelopes through the transport model. The \emph{nonlinear} regime adds the effective
mixing channel without hardware distortion. The \emph{rich} regime includes both. The main
practical objective is a balanced score that rewards detector-limited matched-filter
fidelity while penalizing residual photons, QND defect, leakage, and supra-critical peak
occupancy. Separate information-only and emptying-heavy objectives were retained to expose
local tradeoffs at the representative duration.

Physical-regime optimization used warm starts plus random multistart search. A bounded
Powell refinement stage was used earlier in the linear study, but it was unstable here,
especially for \texttt{procedural\_segments}, so the executed physical sweeps set
\texttt{n\_local = 0}. That change reduced optimizer fragility without materially changing
the frontier conclusion.

\section{Results}

\subsection{Goal 1: Hierarchy of realistic frontiers}
The first result is the duration-dependent hierarchy in Figure~\ref{fig:bounds}. At
\SI{96}{ns}, the detector-limited bound is $0.6664$ and the best rich-model family reaches
$0.6480$, so the physically realistic optimum already lies very close to the measurement
bound. At \SI{240}{ns}, the detector-limited bound increases to $0.9847$ while the best
rich-model family reaches $0.9599$. At \SI{496}{ns}, the rich-model frontier reaches
$0.9982$ against a detector-limited bound of $0.999995$.

The important ranking change is not a collapse of structured control, but a change in
which structure is preferred. The earlier nominal replay favored \texttt{procedural\_segments}
at \SI{96}{ns} and \SI{240}{ns}, and \texttt{nulling\_tail} at \SI{496}{ns}. Once the
rich model is applied, \texttt{ring\_hold} leads the frontier at all three executed
durations. This is already enough to answer the main question: low-dimensional procedural
structure survives, but the freest family is no longer the automatic winner.

\begin{figure}[H]
  \centering
  \includegraphics[width=0.92\linewidth]{../figures/fig1_bound_hierarchy.pdf}
  \caption{Executed hierarchy of ideal, detector-limited, nominal, and rich-model
  frontiers. The realistic conclusion changes qualitatively only after both hardware replay
  and nonlinear-QND penalties are included.}
  \label{fig:bounds}
\end{figure}

\subsection{Goal 2: Representative rich-model family comparison}
Table~\ref{tab:representative} gives the representative \SI{240}{ns} comparison for the
rich-model designs. The best overall practical family is the ring-hold design, which
reaches $F_\eta = 0.9599$ with residual occupancy $0.801$ photons and QND preservation
$0.9921$. The free procedural-segments family is close behind in fidelity
($0.9547$) and slightly lower in residual occupancy ($0.713$), but it suffers a larger
measurement-induced transition probability ($4.66\times10^{-3}$) and lower QND
preservation ($0.9886$).

The smooth families reveal the tradeoff clearly. The Fourier-basis family gives lower
fidelity ($0.9347$) but nearly eliminates excess disturbance
($2.35\times10^{-4}$ measurement-induced transition). The nulling-tail family still empties
the cavity well, but the exact linear cancellation no longer translates into frontier-best
rich-model performance once hardware bandwidth and the disturbance layer are applied. The
balanced higher-dimensional reference is even more informative: it reaches only
$F_\eta = 0.9092$ under the practical objective, which means unrestricted dimensionality is
not the route to the best realistic pulse in this model.

\begin{figure}[H]
  \centering
  \includegraphics[width=\linewidth]{../figures/fig2_rich_frontiers.pdf}
  \caption{Rich-model family frontiers. The best family is ring-hold across the
  executed durations, while square and smooth-square baselines remain clearly suboptimal.}
  \label{fig:frontiers}
\end{figure}

\begin{table}[H]
  \centering
  \small
  \caption{Representative \SI{240}{ns} rich-model comparison. The piecewise reference is
  shown at its balanced operating point, not its information-only optimum.}
  \label{tab:representative}
  \begin{tabular}{@{} l c c c c c @{} }
    \toprule
    Family & $F_\eta$ & $n_{\mathrm{res}}$ & $Q_{\mathrm{QND}}$ & $p_{\mathrm{ind}}$ & $n_{\mathrm{peak}}$ \\
    \midrule
    Square & 0.9201 & 2.366 & 0.9931 & $1.94\times10^{-4}$ & 3.156 \\
    Ring-hold & 0.9599 & 0.801 & 0.9921 & $1.23\times10^{-3}$ & 4.417 \\
    Procedural segments & 0.9547 & 0.713 & 0.9886 & $4.66\times10^{-3}$ & 3.999 \\
    Nulling tail & 0.9216 & 0.887 & 0.9931 & $2.20\times10^{-4}$ & 3.356 \\
    Fourier basis & 0.9347 & 0.673 & 0.9931 & $2.35\times10^{-4}$ & 2.925 \\
    Piecewise ref. (bal.) & 0.9092 & 0.416 & 0.9933 & $4.09\times10^{-5}$ & 3.098 \\
    \bottomrule
  \end{tabular}
  \normalsize
\end{table}

\subsection{Goal 3: Separating hardware transport from nonlinear QND loss}
Figure~\ref{fig:breakdown} isolates how the richer model changes one representative pulse.
For the procedural-segments design at \SI{240}{ns}, the nominal replay achieves
$F_\eta = 0.9628$ and $Q_{\mathrm{QND}} = 0.9933$. Adding only multilevel relaxation does
almost nothing. Adding only hardware replay lowers the fidelity to $0.9552$ with
essentially no extra induced transition. Adding only the nonlinear layer leaves the
fidelity nearly unchanged ($0.9625$) but lowers QND preservation to $0.9898$ through a
measurement-induced transition probability of $3.51\times10^{-3}$. The full rich model
combines both effects, giving $F_\eta = 0.9547$ and
$Q_{\mathrm{QND}} = 0.9886$.

This decomposition matters because it explains why the control ranking changes. Hardware
transport punishes sharp or highly phase-structured envelopes by distorting the waveform.
The nonlinear channel punishes envelopes that achieve their performance by pushing harder
into fast, high-occupancy excursions. A family like ring-hold is smooth enough
to survive the first effect and conservative enough to avoid most of the second.

\begin{figure}[H]
  \centering
  \includegraphics[width=0.92\linewidth]{../figures/fig3_regime_breakdown.pdf}
  \caption{Representative replay breakdown for the procedural-segments design.
  Hardware distortion accounts for most of the fidelity loss, while the nonlinear channel
  accounts for most of the new QND defect.}
  \label{fig:breakdown}
\end{figure}

\begin{figure}[H]
  \centering
  \includegraphics[width=0.92\linewidth]{../figures/fig4_transport_waveforms.pdf}
  \caption{Programmed versus hardware-replayed waveforms for the representative designs.
  Finite bandwidth, timing quantization, and IQ imbalance visibly smooth and skew the most
  aggressive pulse edges.}
  \label{fig:transport}
\end{figure}

\subsection{Goal 4: Tradeoffs, stress, and robustness}
The piecewise reference slice in Figure~\ref{fig:tradeoff} shows why the practical optimum
is not the information-only optimum. At \SI{240}{ns}, the information-optimized reference
reaches $F_\eta = 0.9760$, but only by leaving $4.31$ residual photons and reducing QND
preservation to $0.9802$. Its balanced counterpart drops to $F_\eta = 0.9092$ while
recovering residual occupancy and QND behavior. The lesson is that the remaining detector
gap cannot be interpreted as freely available performance; it is tied to readout-induced
state disturbance and cavity refill cost.

The QND stress test strengthens the same conclusion. At $1.3\times$ the nominal amplitude
probe, \texttt{procedural\_segments} falls to $Q_{\mathrm{QND}} = 0.9618$ with induced
transition probability $3.15\times10^{-2}$. By contrast, \texttt{fourier\_basis} remains
at $Q_{\mathrm{QND}} = 0.9890$ with induced transition $4.33\times10^{-3}$, and the square
baseline remains at $0.9864$. The flexible family is the most sensitive to strong-drive
stress.

Finally, the robustness sweep shows that the practical ordering is not a one-point
artifact. Across the detuning, phase, hardware, and amplitude probes, the nominal, mean,
and worst rich-model fidelities are $(0.9201, 0.9191, 0.8645)$ for square,
$(0.9547, 0.9533, 0.9017)$ for procedural segments,
$(0.9216, 0.9208, 0.8712)$ for nulling tail, and
$(0.9347, 0.9336, 0.8807)$ for the Fourier basis. The procedural family keeps a fidelity
edge under perturbations, but its QND-stress margin is clearly worse than smoother
alternatives.

\begin{figure}[H]
  \centering
  \includegraphics[width=0.92\linewidth]{../figures/fig5_tradeoffs.pdf}
  \caption{Representative tradeoff slice for the higher-dimensional reference family. The
  information-only optimum is not practically acceptable once emptying and QND penalties are
  included.}
  \label{fig:tradeoff}
\end{figure}

\begin{figure}[H]
  \centering
  \includegraphics[width=0.92\linewidth]{../figures/fig6_qnd_stress.pdf}
  \caption{QND stress test under amplified readout drive. The free procedural family is the
  most sensitive to the strong-drive disturbance channel.}
  \label{fig:qndstress}
\end{figure}

\begin{figure}[H]
  \centering
  \includegraphics[width=0.92\linewidth]{../figures/fig7_robustness.pdf}
  \caption{Rich-model robustness summary across hardware and control probes. The fidelity
  advantage of structured optimization survives moderate perturbations, but the family with
  the most freedom is not the most disturbance-robust.}
  \label{fig:robustness}
\end{figure}

\section{Validation}

\subsection{Sanity Checks}
All required sanity checks passed. The analytic nulling-tail construction drives both
conditioned linear cavity amplitudes to machine zero,
$|\alpha_g(T)| = |\alpha_e(T)| = 3.54\times10^{-16}$. The zero-drive replay produces zero
residual photons and only $1.38\times10^{-12}$ induced transition probability, so the
rich-model disturbance channel is not spuriously active in the absence of readout drive.

\subsection{Convergence Analysis}
The representative \SI{240}{ns} rich-model point is stable against both smaller time steps
and larger cavity truncation. The fine-time-step replay changes matched-filter fidelity by
$1.43\times10^{-3}$, and increasing the cavity truncation changes it by only
$4.73\times10^{-5}$. These values satisfy the study's practical convergence target.

\subsection{Literature Comparison}
This study is not a direct reproduction benchmark, so a one-to-one percent-error table is
not appropriate. The comparison to the literature is instead qualitative but still useful.
First, active cavity emptying remains valuable, consistent with driven-depletion results in
the superconducting-qubit readout literature \cite{mcclure2016rapid,walter2017rapid}.
Second, the richer model reproduces the expected control trend that stronger, faster
readout can induce state transitions once one moves beyond the naive dispersive picture,
consistent with the mechanism emphasized in \cite{sank2016measurement}. The present study
therefore agrees with published physical intuition, while remaining honest about the fact
that its nonlinear layer is phenomenological rather than microscopic.

\section{Discussion}
The central lesson is narrower and more useful than either a blanket success or blanket
failure story. Procedural optimization still works, in the sense that low-dimensional,
physics-informed pulse shaping captures most of the useful gain over square readout.
However, the best practical family is not the one with the most unconstrained freedom.
Under the richer model, the winner is a structured and smooth pulse: \texttt{ring\_hold}.

The decomposition in Figure~\ref{fig:breakdown} explains why. The earlier study mostly
ranked families by information extraction plus ordinary cavity emptying. Once hardware
transport is added, envelopes with sharper features lose fidelity through distortion. Once
the nonlinear channel is added, families that rely on rapid occupancy excursions lose QND
performance through induced transitions. The free procedural family can still score very
well on information, but it does so by spending more of the available budget on exactly the
features that the richer model penalizes.

The nulling-tail result is equally instructive. In the linear model, exact cancellation is
a powerful constructive principle. In the richer model, that principle still helps, but it
is no longer the whole problem. A pulse can empty the cavity efficiently and still lose the
frontier because the waveform needed to realize that cancellation is transport-distorted or
because the transient occupancy pattern is no longer QND-friendly.

The high-dimensional reference provides the final control argument. When it is allowed to
optimize only for information, it beats all practical families on $F_\eta$ but leaves too
many photons and too much QND damage. When it is re-scored under the balanced practical
objective, it no longer dominates the structured families. This shows that the realistic
gap is not primarily a missing-optimizer problem; it is a physics-and-constraints problem.

\section{Conclusion}
The question posed at the start of the study can now be answered directly. Physics-informed
procedural readout pulses do remain near-optimal once nonlinear-QND effects and
hardware-realistic transport are included, but only in the weaker sense that structured
pulse shaping remains valuable. They do \emph{not} remain near-optimal in the stronger
sense that the freest procedural family is still the best practical default.

Within the present rich-model hierarchy, the recommended family is
\texttt{ring\_hold}. It leads the executed frontier at \SIlist{96;240;496}{ns}, preserves
high QND quality, and is materially less fragile than the free procedural family under
strong-drive stress. \texttt{procedural\_segments} remains a useful specialist option when
raw information fidelity is the dominant objective and the user is willing to accept a
tighter QND budget. \texttt{nulling\_tail} remains valuable when cavity emptying is the
main concern, but it is no longer the overall frontier leader once transport and nonlinear
backaction are both present.

The study goals were therefore met. The replay stack now includes explicit hardware
transport and family-dependent non-QND disturbance, the practical frontier was re-ranked,
the family comparison was validated, and a concrete control recommendation now exists.

\section{Limitations and Future Work}

\subsection{Known Limitations}
The strongest limitation is model fidelity. The effective mixing channel is intentionally
phenomenological, constructed to expose the control consequences of strong-drive QND
breakdown without claiming a microscopic derivation. That is sufficient for ranking trends,
but it is not sufficient for quantitative hardware signoff.

The second limitation is coverage. The stable saved frontier only spans
\SIlist{96;240;496}{ns}, because the rich-model optimization was significantly more
expensive and more fragile than the nominal study. That grid is enough to answer the main
question, but it is not enough to map every duration crossover precisely.

The third limitation is the replay formalism. The public \texttt{cqed\_sim} workflow still
uses deterministic Lindblad replay rather than a stochastic measurement unraveling with
time-dependent collapse structure. Repeated-readout quality is therefore quantified through
matched replay metrics rather than trajectory-conditioned state update.

\subsection{Suggested Improvements}
\begin{itemize}
  \item \textbf{[P1 | HIGH]} Calibrate the effective mixing layer against a more microscopic
  beyond-dispersive or experimentally informed model. The present ranking is plausible, but
  its absolute transition probabilities should not yet be treated as device-calibrated.
  \item \textbf{[P1 | MEDIUM]} Run a denser duration sweep in the \SIrange{200}{350}{ns}
  window. This is where the ranking changes from nominal \texttt{procedural\_segments}
  dominance to rich-model \texttt{ring\_hold} dominance.
  \item \textbf{[P2 | MEDIUM]} Introduce a constrained hybrid family built around
  \texttt{ring\_hold} with a small number of additional smooth degrees of freedom. The goal
  is to test whether the remaining detector gap can be closed without reproducing the QND
  fragility of free procedural segments.
  \item \textbf{[P2 | HIGH]} Extend the study to a full three-mode replay or a more explicit
  readout-plus-storage validation layer once compute permits. The present conclusion rests
  on the two-mode-plus-mixing hierarchy.
  \item \textbf{[P3 | MEDIUM]} Add trajectory-based repeated-readout diagnostics when the
  simulator supports them. That would better separate detector noise, state disturbance,
  and readout-conditioned collapse.
\end{itemize}

\subsection{Open Questions}
The most important unresolved question is whether a lightly augmented smooth family can
match the information performance of \texttt{procedural\_segments} without inheriting its
QND-stress penalty. A second open question is whether the nulling-tail family loses the
frontier mainly because finite bandwidth destroys exact cancellation or because the richer
model rewards smooth occupancy trajectories more than exact emptying per se. Finally, the
relative roles of occupancy-driven versus slew-driven disturbance remain only partially
separated in the current phenomenological model.

\bibliographystyle{unsrtnat}
\bibliography{references}

\end{document}